\pdfbookmark[1]{Abstract}{Abstract}
\chapter*{Abstract}

Reconnaissance, not a result. What the Millennium Prize problems ask, in the words of their own
official statements. What became of one of them in September 2026. Which of the rest this project's
apparatus could be aimed at. And the work nobody here has done. No problem gets attacked in these
pages, and nobody claims priority over anything.

Seven documents make up the corpus. Five are Clay Mathematics Institute problem descriptions,
covering P versus NP, the Hodge conjecture, the Riemann hypothesis, quantum Yang--Mills theory, and
the Birch and Swinnerton-Dyer conjecture. The other two are Riemann: his 1859 memoir in the Wilkins
translation, and four photographed pages of his handwriting. Somebody opened each one. The
photographs turn out to hold a draft, not the published memoir, and no text can be pulled from them
at all, though the filename says otherwise. Seed a citation off that filename and the registry
inherits the error. The corpus chapter exists for that reason.

Navier--Stokes is missing from the corpus because the Clay index of unsolved problems no longer
carries it. OpenAI published a finite-time blowup proof for the forced three-dimensional
incompressible equations on 8 September 2026, with a Lean~4 formalization beside it. Fefferman asks
for a proof of any one of four statements. Two of those four put a smooth external force inside the
statement. Forcing sits in the problem as posed, then. Twice this work insisted otherwise before
anyone opened the document, and the last chapter carries both retractions along with whatever
killed them.

Six points went into checking that formalization against the problem statement. Viscosity strictly
positive, with nothing taken to a vanishing limit. Decay conditions on the data and on the force,
neither relaxed. A force smooth across the whole closed half-line. A conclusion quantified as a
negated existential. Kinetic energy held bounded for all time, not merely at the start. And the
periodic pressure condition that arrives in the errata and not in the body. All six match.

Whose statement is it, though. Not the authors' own. They took it from a Lean formalization the
Formal Conjectures project had already published, kept it unmodified, and left its deliberate
placeholder gaps in place, while a separate module discharges the identical statement and prints out
what axioms it leaned on. A proof assistant checks a deduction and never a translation. That is this
project's standing objection to machine-checked proof, and the architecture above is built to answer
it.

The gaps matter as much, so they sit in the sentences that carry the claims and not in a footnote.
Nobody built the Lean development. Nobody machine-checked anything here. Source text is what got
read. Of the paper, the statements and the bridging argument, nothing beyond. The proof tree stayed
shut. Nobody looked into the priority dispute running through this field, and no line here should be
read as taking a side in it.

Which problem can the apparatus reach? Birch and Swinnerton-Dyer, on the project architect's
recommendation, which this book records without adopting. Their reasoning: the refined conjecture is
an exact numerical identity, so verifying it comes down to whether a computed real number lands on
an integer. A precision engine answers questions shaped like that one. Riemann goes in the same
chapter as the trap. There the apparatus matches the subject and misses the obstruction, zeros on
the critical line have been computed far past anything reachable here, and counting more of them
bears on no proof. One null offered for the Birch and Swinnerton-Dyer measurement got withdrawn by
its own author, who had derived the bar instead of drawing it; the drawn replacement stands in its
place.

Nothing here claims a new method. Nothing here attributes one either. The literature behind these
problems has gone unread in this work, and until somebody reads it that stays missing work, not a
hole in anyone else's.
\cleardoublepage
